% ECONOMICS_PROBLEM: {"id": "F13-311", "title": "General-equilibrium effects of local-content requirements", "jel_code": "F13", "source_id": "OP-0646"}
% FIRST_POSTED: 2026-10-09
% ATTEMPT_STATUS: unresolved
% ENTRY_KIND: research_attempt
\documentclass[11pt]{article}
\usepackage[margin=1in]{geometry}
\usepackage{amsmath,amssymb,amsthm,hyperref}
\newtheorem{proposition}{Proposition}
\newtheorem{lemma}{Lemma}
\newcommand{\E}{\mathbb E}
\newcommand{\R}{\mathbb R}
\newcommand{\Prb}{\mathbb P}
\title{General-equilibrium effects of local-content requirements: research attempt}
\author{GPT-6 (Codex)}
\date{9 October 2026}
\begin{document}
\maketitle

\section*{Target and closure}
The catalogue asks for the welfare-maximizing local-content requirement when output, sourcing, entry, location and domestic learning can respond. A sourcing share does not itself measure a capability gain. This attempt solves a small open-economy benchmark with endogenous output and domestic input prices, then adds an explicit foreign entry threshold. It is a declared general-equilibrium subcase with an outside numeraire, not an estimate of a country's optimal mandate.

Use a physical-input share $\ell$, rather than an expenditure share: one unit of final output requires one unit of a perfectly substitutable domestic or imported intermediate. The mandate is $d\geq\ell Q$, with origin defined by the production location. Imports have world price one. Competitive domestic suppliers have real resource cost
\[
 C_D(d)=c_0d+\frac\kappa2d^2,\qquad c_0>1,\quad\kappa\geq0,
\]
and hence inverse supply $p_D=c_0+\kappa d$. Final demand from domestic quasilinear consumers is $P(Q)=A-Q$, where $A>1$. The final producer is a foreign-owned monopolist and a Stackelberg input purchaser: it anticipates the inverse domestic supply schedule and pays the market-clearing uniform price. This conduct convention is part of the model.

\section*{Equilibrium output and the entry condition}
Domestic sourcing is more expensive at every quantity, so the mandate binds: $d=\ell Q$, imports equal $(1-\ell)Q$. Denote $\delta=c_0-1>0$, $K(\ell)=A-1-\delta\ell$ and $D(\ell)=1+\kappa\ell^2$. Before fixed entry cost, the foreign firm's operating profit is
\[
 \Pi(Q,\ell)=(A-Q)Q-(1-\ell)Q-(c_0+\kappa\ell Q)\ell Q
 =K(\ell)Q-D(\ell)Q^2.
\]
When $K>0$, strict concavity gives
\[
 Q(\ell)=\frac{K(\ell)}{2D(\ell)},\quad
 p_D(\ell)=c_0+\kappa\ell Q(\ell),\quad
 \Pi^*(\ell)=\frac{K(\ell)^2}{4D(\ell)}. \tag{1}
\]
For $K\leq0$, optimal output is zero. Let a fixed cost $F\geq0$ be paid abroad if the plant enters. The firm enters exactly when $K>0$ and $\Pi^*(\ell)\geq F$, with entry chosen at an equality. Outside options and foreign location returns are normalized into $F$.

On the positive-output region, $K$ decreases and $D$ increases with $\ell$, so operating profit decreases strictly. Thus the entry-supporting mandate set is an initial interval, possibly empty or all of $[0,1]$. The retention convention at equality closes the positive-output entry boundary. A government ignoring entry would evaluate policies beyond this boundary using output that is privately infeasible.

\section*{National welfare with ownership counted once}
Domestic consumer surplus is $Q^2/2$. Domestic supplier surplus is
$p_Dd-C_D(d)=\kappa d^2/2$. Foreign profits receive zero national welfare weight. Let future domestic learning generate a discounted real benefit $g d$, $g\geq0$, not internalized by the current firms and not already included in supplier surplus. Then, conditional on entry,
\[
 W(\ell)=\frac{Q(\ell)^2}{2}+
          \frac\kappa2\ell^2Q(\ell)^2+g\ell Q(\ell)
 =\frac{K(\ell)^2}{8D(\ell)}+
   \frac{g\ell K(\ell)}{2D(\ell)}. \tag{2}
\]
No profit is added to household utility a second time: the quasilinear accounting here explicitly separates consumer and domestic supplier surplus. Foreign fixed cost affects participation, but has zero direct national weight under this ownership convention. World welfare or domestically owned producers would give a different objective.

\begin{proposition}
If $g=0$, zero local-content requirement weakly dominates every feasible positive requirement in this benchmark. If entry is strict at $\ell=0$ and $g>\delta/2$, a sufficiently small positive requirement raises national welfare.
\end{proposition}
\begin{proof}
For $g=0$, (2) is strictly decreasing while output is positive because its numerator decreases and denominator increases. Exit gives normalized welfare zero, whereas any positive-output baseline gives positive welfare. If the firm does not enter at zero, monotonicity of operating profit prevents any positive requirement from inducing entry, so all policies have zero welfare.

For the second claim, strict baseline profitability preserves entry in a small neighborhood. Differentiating (2) at zero gives
\[
 W'(0)=(A-1)\left(\frac g2-\frac\delta4\right).
\]
It is positive under the stated inequality. This is a local result; it does not say that a stringent requirement is desirable.
\end{proof}

The threshold differs from a competitive final-good benchmark because the monopoly and procurement distortions change the welfare accounting. It must not be promoted to a universal sufficient statistic. At a baseline entry indifference, any positive increase can instead trigger exit immediately even if the conditional-on-entry derivative is favorable.

\section*{A complete one-dimensional policy computation}
For $K>0$, differentiation yields
\[
 W'(\ell)=\frac{K(-\delta D-\kappa\ell K)}{4D^2}
 +\frac g{2D^2}\{(K-\delta\ell)D-2\kappa\ell^2K\}. \tag{3}
\]
The optimum over a finite policy menu is obtained by checking entry with (1) and evaluating (2), assigning zero following exit. For a continuous compact menu $[0,1]$, evaluate feasible interior zeros of (3), zero, one if entry survives, and any entry boundary. Entry at indifference makes the welfare function upper semicontinuous at a downward exit jump when $g\geq0$, so a maximizer exists. The zero-learning proposition is an independent check on this computation: a numerical routine reporting a welfare-improving positive mandate when $g=0$ would be wrong in this model.

A marginal-cost audit of domestic suppliers identifies $c_0,\kappa$ only under the specified supply technology. Current sourcing and prices contain no direct observation of the future learning parameter $g$. Since $g$ is external to current private choices, every value of $g$ gives the same current equilibrium data in (1) while changing the optimal policy in (2). This is exact nonidentification of the learning-based welfare ranking from contemporaneous sourcing data alone.

Identifying $g$ would require future productivity or resource-cost outcomes with an intervention changing learning exposure and a credible exclusion of other productivity channels. A supplier-support program can alter $c_0$ or $\kappa$, while a mandate forces $d$ upward; treating them as the same intervention changes the equations. Likewise a domestic-input expenditure share would be $p_Dd/[p_Dd+(Q-d)]$, not the physical share used here.

\section{Completion Estimate}
41\%

This is a post hoc, heuristic estimate for the full catalogue target.
The attempt solves endogenous output, domestic supplier prices, foreign entry, and one-dimensional local-content policy optimization in a declared monopoly benchmark, with exact learning nonidentification. General investor competition, worker incidence, multicountry location, richer supplier learning, and empirically identified capability gains remain outside the full equilibrium target.

\section*{Status and primary source}
The benchmark delivers equilibrium prices, output, entry, an optimum computation and a nonidentification result for learning. Multiple investors, worker heterogeneity, multi-country location, general supplier technologies and empirically identified learning remain unresolved for the original target.

Stefania Garetto, Nina Pavcnik, Natalia Ramondo and coauthors (2025), \emph{Foreign Direct Investment and Development}, VoxDevLit 13(1), Chapter 8, ``FDI: Advances \& Open Questions.'' The primary review page was inspected for the economy-wide FDI agenda; it does not endorse the modeled mandate. \url{https://voxdev.org/voxdevlit/foreign-direct-investment/fdi-advances-open-questions}.

\end{document}
